\chapter{Metrics}
\label{ch:metrics}

Venue processes expose metrics that describe message handling, queueing, matching-engine
activity, and resource use. Prometheus collects these measurements, and the Python metrics
visualiser and Grafana dashboards make them available for investigation. The measurements
describe only the events and values recorded by the instrumented code; their names and
limitations determine what can be concluded from them.

\section{Collection and configuration}

Metrics are configured separately for each process. When metrics are enabled, a process
provides an HTTP endpoint that Prometheus scrapes. The endpoint does not require
authentication and does not encrypt its traffic, so its listening address must be chosen
with the surrounding network and access controls in mind.

The metrics configuration identifies the application, the process instance, and the
address and port used by the endpoint. Omitting the metrics configuration leaves
collection disabled. When disabled, the framework uses no-op metric handles. Calls that
record a metric still evaluate their arguments, but the values are not exported.

The supplied development configuration enables venue metrics and configures Prometheus
to scrape every five seconds. The supplied test, pre-production, and production
configurations disable venue metrics. The development Prometheus configuration retains
samples for 90 days. The development environment can also collect host-level measurements,
which can be viewed alongside the venue's process metrics.

\section{Metric names and types}

Each metric key has three or four dot-separated parts:
\begin{itemize}
  \item the application name;
  \item the process instance name;
  \item an optional scope, such as a thread or subsystem;
  \item the metric name.
\end{itemize}

The application, process instance, and optional scope become Prometheus labels. The final
part is the exported metric name. For example, the key
\texttt{pubsub.fix\_order\_gateway\_a.gateway\_thread.order\_round\_trip\_nanoseconds}
becomes the metric \texttt{order\_round\_trip\_nanoseconds} with the labels
\texttt{application="pubsub"}, \texttt{component="fix\_order\_gateway\_a"}, and
\texttt{scope="gateway\_thread"}. The process label identifies a particular process
instance, rather than a general component type.

The framework exposes its reactor measurements for every process with metrics enabled.
Application measurements are opt-in: the sequencer, matching engine, and order gateways
configure a scope for their application threads. A thread's scope distinguishes its
measurements from those of other threads. This is also important for histograms, which
must be recorded by one thread only.

Counters record values that accumulate, such as the number of orders processed. Gauges
record values that can rise or fall, such as queue depth or the number of objects currently
in a pool. Histograms record observations in configured buckets, along with a count and
sum. Prometheus adds a cumulative count for each bucket boundary, as well as a bucket for
all observations. Most latency measurements are recorded in nanoseconds.

Some memory-pool measurements are cumulative counts but are exported as gauges. Queries
and dashboards must account for that distinction rather than assuming that a metric's
meaning can always be inferred from its Prometheus type.

\section{Measurements recorded}

The framework records measurements at several points in message processing. These
measurements show where work is performed and how long it takes, but they do not all
describe the same interval.

\paragraph{Network and framework messages.}
\texttt{framework\_pdu\_messages\_total} counts protocol data units delivered to
application threads. It does not count every PDU read from the network: a PDU rejected
before delivery is not included. The count is recorded per thread, so a process-wide
total must be formed by summing the thread series.

\paragraph{Inter-thread queues and reactor work.}
\texttt{itc\_queue\_depth} records how many messages remain behind the message being
dispatched. \texttt{itc\_queue\_latency\_nanoseconds} measures from enqueueing until
dispatch. That interval includes both time waiting in the queue and the notification and
wake-up needed to dispatch it. Reading the queue depth alongside the latency can help
distinguish time spent waiting for a thread to run from time spent behind other queued
messages.

The reactor records separate histograms for command delay, sending, receiving, and its
own work cycle. \texttt{reactor\_command\_latency\_nanoseconds} measures from when an
application thread queues a command until the reactor takes it up.
\texttt{reactor\_send\_path\_nanoseconds} measures the reactor's work to turn a send
request into bytes on a socket, and
\texttt{reactor\_receive\_path\_nanoseconds} measures its work to turn readable bytes
into a message on an application thread's queue.
\texttt{reactor\_lap\_nanoseconds} measures the time between successive checks for work.
These are different intervals and should not be treated as interchangeable measures of
network or application latency.

\paragraph{Write-ahead log.}
\texttt{wal\_append\_nanoseconds} records the time spent committing a record to the
write-ahead log. It measures the append operation, not the full time for an order to reach
the matching engine or a client.

\paragraph{Order processing and replication.}
\texttt{orders\_processed\_total} increases when the matching engine inserts a new order
into its book. It excludes replayed orders, orders discarded by a follower, duplicate
orders rejected by the matching engine, and cancellations. It is not a count of every
order-related event handled by the venue.
\texttt{follower\_report\_differences\_total} counts differences between reports produced
by the follower and the leader while applying the same write-ahead log. A non-zero value
indicates that the two matching engines did not produce identical reports.

\paragraph{Gateway activity and latency.}
Gateway counters record activity such as orders received, execution reports sent or
dropped, orders refused, sessions established or lost, and orders rejected because a
session has reached its rate limit. These counts describe the gateway's processing and
session state; they are not interchangeable with counts recorded by the sequencer or
matching engine.

\texttt{order\_round\_trip\_nanoseconds} measures from when an order arrives at a gateway
until that gateway starts sending the ExecutionReport that acknowledges a new order. It
does not measure when the client receives the report. The separate
\texttt{order\_ingress\_to\_forward\_nanoseconds} histogram ends when the gateway forwards
the order to the sequencer. It measures gateway work, including parsing and validation,
and is the more direct comparison of the FIX and binary gateways. Their round-trip
measurements also include the work shared by the sequencer, matching engine, and report
path.

\texttt{order\_path\_elapsed\_nanoseconds} records how long an order has been inside the
venue at several points on its path. The scope identifies the point, such as
\texttt{order\_in}, \texttt{order\_out}, \texttt{er\_in}, or \texttt{er\_out}. These
readings start from the same gateway ingress time. Subtracting readings at consecutive
points gives the time spent between those points; a reading by itself describes how far
along the path the order has travelled, not the duration of a separate operation.

\paragraph{Order-book and memory-pool state.}
Order-book gauges report the number of entries, occupied slots, whether the book is being
moved into a larger table, and the largest table allocation made. Memory-pool gauges
report values such as objects allocated or available and bytes in use or reserved. Pool
measurements are scoped to an individual pool. Some pool gauges, including allocation
and expansion counts, are cumulative despite being exported as gauges; queries must
account for that when calculating changes over time.

\section{Viewing the measurements}

The Python visualiser in \texttt{scripts/pubsub\_metrics.py} queries Prometheus and can
list the applications, process instances, and metrics currently available. It can display
counters, gauges, and histogram data for a selected process. It also supports saved
snapshots for later inspection and synthetic data for demonstrating its displays without
using live measurements.

For latency histograms, the visualiser can show estimated percentiles and latency bands.
The bands report the 50th, 90th, and 99th percentiles together with the number of
observations, so latency can be considered alongside the amount of traffic measured.
Optional ceiling counts are available only when the histogram has a bucket boundary at
the requested ceiling. The visualiser can compare all order gateways for a selected
histogram, save a snapshot for rendering on another machine, and display synthetic
demonstration data without querying Prometheus.

The dashboard generator in \texttt{scripts/make\_grafana\_dashboard.py} creates panels
for latency, throughput, queueing, order-book state, and memory-pool pressure. Its
protocol comparisons include offered rates. Comparing latency without considering the
amount of offered work can make differences in workload appear to be differences in
protocol behaviour.

\section{Interpreting results}

Histogram percentiles are estimates based on the configured bucket boundaries. A
percentile that falls within a bucket is estimated from that bucket; it is not a record
of an individual observation at that value. Histograms do not provide the true minimum or
maximum, and they do not preserve individual outlier observations. A ceiling count is
exact only when the ceiling matches a configured bucket boundary.

The observation count matters when reading a percentile. A percentile based on very few
observations may not represent the wider workload. A period with no observations has no
latency percentile to report; it does not mean that latency was zero.

Measurements should be interpreted according to the point where they are recorded. For
example, a count of messages delivered to an application thread cannot be used as the
count of messages read from a network connection. Likewise, the matching engine's
processed-order count cannot be used as a count of all orders received. A metric is
evidence about its defined measurement point, not about events outside that point.
